\exercisesfor{2}

\begin{enumerate}
\def\labelenumi{\arabic{enumi}.}
\item
  We adopt the notation T, J, \(T_{n}\) of Definition 1 and no. 6. Let t be a homogeneous tensor of T of order p and \(\sigma\) a permutation of \(\{1, 2, \ldots, p\}\) . Then \(t - \sigma t \in J + T_{p-1}\) . (Reduce it to the case where \(\sigma\) is a transposition of two consecutive integers.)
\item
  Suppose that \(\mathbf{K}\) is a field. Let \(\mathfrak{g}\) be a Lie \(\mathbf{K}\) -algebra and \(\mathbf{U}\) its enveloping algebra.
\end{enumerate}

\begin{enumerate}
\def\labelenumi{(\alph{enumi})}
\item
  Let \(u, v\) be in U. If \(u\) is of filtration \(> n\) , and \(v\) of filtration \(> p\) , then \(uv\) is of filtration \(> n + p\) (use Theorem 1).
\item
  Deduce from (a) that the only invertible elements of U are the scalars.
\item
  Deduce from (b) that the Jacobson radical of U is zero.
\end{enumerate}

\begin{enumerate}
\def\labelenumi{\arabic{enumi}.}
\setcounter{enumi}{2}
\item
  Suppose that K is a field. Let g be a Lie K-algebra and U its enveloping algebra. The left regular representation of U corresponds to a representation \(\rho\) of g on the space U. Show that the set \(U_{+}\) of elements of U with no constant term is stable under \(\rho\) and that \(U_{+}\) has no supplement in U stable under \(\rho\) . In particular, \(\rho\) is not semi-simple. (This will be considered again in Theorem 2 of § 6.)
\item
  Suppose that K is a field.
\end{enumerate}

\begin{enumerate}
\def\labelenumi{(\alph{enumi})}
\item
  Verify that there exists a 3-dimensional Lie algebra g with basis \((x, y, z)\) such that \([x, y] = z\) , \([x, z] = [y, z] = 0\) . Let U be the enveloping algebra of g. Show that the centre of U is the subalgebra generated by 1 and z. (Consider, for all \(\xi \in g\) , the derivation of the symmetric algebra S of g which extends ad, \(\xi\) and look for the elements of S annihilated by these derivations; then apply the final remark of no. 8.)
\item
  Verify that there exists a 3-dimensional Lie algebra g with basis \((x, y, z)\) such that \([x, y] = y\) , \([x, z] = z\) , \([y, z] = 0\) . Show that the centre of the enveloping algebra of g reduces to the scalars. (Same method.)
\end{enumerate}

\begin{enumerate}
\def\labelenumi{\arabic{enumi}.}
\setcounter{enumi}{4}
\item
  Suppose that \(\mathbf{K}\) is a field of characteristic \(p > 0\) ( \(p\) prime). Let \(\mathfrak{g}\) be a Lie algebra over \(\mathbf{K}\) with a basis, canonically identified with a submodule of its enveloping algebra U. For an endomorphism \(\phi\) of the K-module \(\mathfrak{g}\) to be a \(p\) -mapping, it is necessary and sufficient that \(x \mapsto x^p - \phi(x)\) be a semi-linear mapping (for \(\lambda \mapsto \lambda^p)\) of \(\mathfrak{g}\) into the centre of U. Deduce that if \((b_{\lambda})\) is a basis of \(\mathfrak{g}\) , for there to exist a \(p\) -mapping on \(\mathfrak{g}\) , it is necessary and sufficient that, for all \(\lambda\) , there exist \(c_{\lambda} \in \mathfrak{g}\) such that \((\operatorname{ad} b_{\lambda})^p = \operatorname{ad} c_{\lambda}\) ; if this is so, there then exists one and only one \(p\) -mapping \(x \mapsto x^{[p]}\) such that \(b_{\lambda}^{[p]} = c_{\lambda}\) for all \(\lambda\) .
\item
  Let g be a Lie p-algebra over a ring K such that \(pK = \{0\}\) (p prime), U its enveloping algebra and \(\sigma\) the canonical mapping of g into U. Let J be the two-sided ideal of U generated by the elements \((\sigma(x))^{p} - \sigma(x^{[p]})\) where x runs through g. The associative algebra \(\tilde{U} = U/J\) is called the restricted enveloping algebra of g. The mapping \(\sigma\) defines when passing to the quotient a mapping \(\tilde{\sigma}\) (called canonical) of g into \(\tilde{U}\) , which is a p-homomorphism (when \(\tilde{U}\) is considered as a Lie p-algebra).
\end{enumerate}

\begin{enumerate}
\def\labelenumi{(\alph{enumi})}
\item
  The algebra \(\tilde{\mathbf{U}}\) and the mapping \(\tilde{\sigma}\) are solutions of a universal mapping problem: for every \(p\) -homomorphism \(f\) of \(\mathfrak{g}\) into an associative algebra \(\mathbf{B}\) over \(\mathbf{K}\) (considered as a Lie \(p\) -algebra), there exists one and only one K-homomorphism \(f'\) of \(\tilde{\mathbf{U}}\) into \(\mathbf{B}\) (with their associative algebra structures) such that \(f = f' \circ \tilde{\sigma}\) .
\item
  Show that, if \((b_{\lambda})_{\lambda \in \Lambda}\) is a basis of \(\mathfrak{g}\) (where \(\Lambda\) is totally ordered), \(\tilde{\sigma}\) is injective and that, if \(x\) and \(\tilde{\sigma}(x)\) are identified for \(x\in \mathfrak{g}\) , the elements \(\Pi_{\lambda}b_{\lambda}^{\nu_{\lambda}}\) (where the \(\lambda\) are in increasing order, the \(\nu_{\lambda}\) are zero except for a finite number and \(0\leqslant \nu_{\lambda} < p\) for all \(\lambda\) ) form a basis of \(\tilde{\mathbf{U}}\) . (Canonically identifying \(\mathfrak{g}\) with a submodule of \(\mathrm{U}\) , we write \(\phi (x) = x^{p} - x^{[p]}\) for all \(x\in g\) . For every composite index \(\alpha = (\alpha_{\lambda})\in \mathbf{N}^{(\Lambda)}\) , let \(\alpha_{\lambda} = \beta_{\lambda} + p\gamma_{\lambda}\) with \(0\leqslant \beta_{\lambda} < p\) and let
\end{enumerate}

\[
\mathrm{T} _ {\alpha} = (\Pi_ {\lambda} b _ {\lambda} ^ {\beta_ {\lambda}}) (\Pi_ {\lambda} (\phi (b _ {\lambda})) ^ {\gamma_ {\lambda}}).
\]

Show that the \(T_{\alpha}\) form a basis of U and the \(T_{\alpha}\) such that \(\gamma = (\gamma_{\lambda}) \neq 0\) form a basis of J; observe that the \(\phi(b_{\lambda})\) belong to the centre of U.)

\begin{enumerate}
\def\labelenumi{\arabic{enumi}.}
\setcounter{enumi}{6}
\tightlist
\item
  Let g be a Lie p-algebra over a ring K such that \(pK = \{0\}\) (p prime). A derivation D of g is called a p-derivation if \(\mathrm{D}(x^{p}) = (\mathrm{ad} x)^{p-1}\) . Dx for all \(x \in g\) . Every inner derivation is a p-derivation.
\end{enumerate}

\begin{enumerate}
\def\labelenumi{(\alph{enumi})}
\item
  If \(\mathbf{L}\) is an associative K-algebra, every derivation of \(\mathbf{L}\) is a \(p\) -derivation when \(\mathbf{L}\) is considered as a Lie \(p\) -algebra. (Use formula (2) of Exercise 19 of § 1.)
\item
  Suppose that g has a basis. For a derivation of g to be a p-derivation, it is necessary and sufficient that it can be extended to a derivation of the restricted enveloping algebra of g. Deduce that the p-derivations of g form a Lie p-subalgebra of the p-algebra of derivations of g.
\item
  If D is a \(p\) -derivation of \(\mathfrak{g}\) , the kernel of D is a \(p\) -subalgebra of \(\mathfrak{g}\) .
\item
  For every derivation D of g, \(\mathrm{D}(x^{p}) - (\mathrm{ad} x)^{p-1}\) . \(\bar{D}x\) belongs to the centre of g for all \(x \in g\) (use formula (2) of Exercise 19 of §1, applied to \(\mathrm{L} = \mathcal{L}(\mathfrak{g})\) ).
\end{enumerate}

\begin{enumerate}
\def\labelenumi{\arabic{enumi}.}
\setcounter{enumi}{7}
\item
  Show that Theorem 1 remains valid if the module \(\mathfrak{g}\) is a direct sum of monogenous modules. (Replace the module P in the proof by the symmetric algebra of g.)
\item
  Let \(k\) be the field with two elements, \(\mathbf{V}\) the vector space \(k^3\) , \((x_1, x_2, x_3)\) its canonical basis and \(\mathbf{K}\) the exterior algebra of \(\mathbf{V}\) , which is an 8-dimensional commutative algebra over \(k\) . Let \(\mathfrak{g}\) be the Lie K-algebra admitting the basis \((e_1, e_2, e_3, e_{12}, e_{13}, e_{23})\) such that
\end{enumerate}

\[
\left[ e _ {1}, e _ {2} \right] = \left[ e _ {2}, e _ {1} \right] = e _ {1 2} \quad \left[ e _ {1}, e _ {3} \right] = \left[ e _ {3}, e _ {1} \right] = e _ {1 3}, \quad \left[ e _ {2}, e _ {3} \right] = \left[ e _ {3}, e _ {2} \right] = e _ {2 3}
\]

and the other brackets are zero. Let \(\mathfrak{h}\) be the ideal of \(\mathfrak{g}\) generated by \(u = x_{1}e_{1} + x_{2}e_{2} + x_{3}e_{3}\) . As a module, \(\mathfrak{h}\) is generated by \(u\) , \([u, e_1] = x_2e_{12} + x_3e_{13}\) , \([u, e_2] = x_1e_{12} + x_3e_{23}\) , \([u, e_3] = x_1e_{13} + x_2e_{23}\) . Let

\[
v = x _ {1} x _ {2} e _ {1 2} + x _ {1} x _ {3} e _ {1 3} + x _ {2} x _ {3} e _ {2 3}.
\]

\[
\phi \left(e _ {1}\right) = \phi \left(e _ {2}\right) = \phi \left(e _ {3}\right) = 0, \phi \left(e _ {1 2}\right) = x _ {3}, \phi \left(e _ {1 3}\right) = x _ {2}, \phi \left(e _ {2 3}\right) = x _ {1}.
\]

\begin{enumerate}
\def\labelenumi{(\alph{enumi})}
\setcounter{enumi}{1}
\item
  Let \(f\) be a linear mapping of \(\mathfrak{g}\) into an associative K-algebra such that \(f([x,y]) = f(x)f(y) - f(y)f(x)\) for all \(x, y\) in \(\mathfrak{g}\) . Show that \(f(v) = f(u)^2\) .
\item
  Deduce from (a) and (b) that the canonical mapping of the Lie algebra \(\mathfrak{g} / \mathfrak{h}\) into its enveloping algebra is not injective.
\end{enumerate}

\begin{enumerate}
\def\labelenumi{\arabic{enumi}.}
\setcounter{enumi}{9}
\tightlist
\item
  Let \(\mathfrak{g}\) be a finite-dimensional Lie algebra over a field, U its enveloping algebra and \(\mathbf{U}_n\) the set of elements of U of filtration \(\leqslant n\) .
\end{enumerate}

\begin{enumerate}
\def\labelenumi{(\alph{enumi})}
\item
  Let \(x \in \mathrm{U}_n, y \in \mathrm{U}_n\) be two non-zero elements. Show that there exists \(u \in \mathrm{U}, v \in \mathrm{U}\) such that \(ux = vy\) . (Compare \(\dim(\mathrm{U}_p x) = \dim \mathrm{U}_p\) , \(\dim(\mathrm{U}_p y) = \dim \mathrm{U}_p\) and \(\dim \mathrm{U}_{p+n}\) . Deduce that \(\mathrm{U}_p x \cap \mathrm{U}_p y \neq \{0\}\) for \(p\) sufficiently large.)
\item
  Show that U admits a field of left quotients (Algebra, Chapter I, § 9, Exercise 8) which is at the same time a field of right quotients.
\end{enumerate}

